\documentclass[10pt,a4paper]{amsart}
\usepackage[margin=19mm]{geometry}
\usepackage{amsmath,amssymb,mathtools}
\usepackage[scr=rsfs,cal=euler]{mathalfa}
\usepackage{xfrac}
\usepackage[unicode,hidelinks,bookmarks=false]{hyperref}
\newcommand{\ie}{i.e.\ }
\newcommand{\cf}{cf.\ }
\newcommand{\resp}{respectively\ }
\newcommand{\Subsection}{\subsection}
\providecommand{\nobreakdash}{\nobreak\hbox{-}}
\setlength{\emergencystretch}{2em}
\multlinegap0pt
\DeclareMathSymbol{\leqslant}{\mathalpha}{AMSa}{"36}
\DeclareMathSymbol{\geqslant}{\mathalpha}{AMSa}{"3E}
\DeclareMathSymbol{\eset}{\mathalpha}{AMSb}{"3F}
\renewcommand{\leq}{\;\leqslant\;}
\renewcommand{\geq}{\;\geqslant\;}
%\newcommand{\dd}{\,\text{\rm d}}




\newcommand{\la}{\label}
\newcommand{\Om}{\Omega}


\newcommand{\e}{\varepsilon}
\newcommand{\h}{\frac{1}{2}}

\newcommand{\rhoN}{\rho^{\pi}}
\newcommand{\si}{\sigma}
\newcommand{\tc}{\, |\, }

\newcommand{\scalar}[2]{\langle #1, #2\rangle}



\DeclareMathOperator*{\union}{\bigcup}
\DeclareMathOperator*{\inter}{\bigcap}
\DeclareMathOperator{\var}{Var}
\newcommand{\tv}{\mathrm{TV}}
\newcommand{\ent}{\mathrm{Ent}}
\newcommand{\gap}{\mathrm{gap}}
\newcommand{\rBe}{\mathrm{Be}}
\newcommand{\IND}{\mathbf{1}}
\newcommand{\ind}{\mathbf{1}}

\newcommand{\cA}{\mathcal{A}}
\newcommand{\cB}{\mathcal{B}}
\newcommand{\cE}{\mathcal{E}}
\newcommand{\cL}{\mathcal{L}}
\newcommand{\cP}{\mathcal{P}}
\newcommand{\cS}{\mathcal{S}}
\newcommand{\cT}{\mathcal{T}}
\newcommand{\cU}{\mathcal{U}}
\newcommand{\cV}{\mathcal{V}}
\newcommand{\bbC}{\mathbb{C}}
\newcommand{\bbE}{\mathbb{E}}

\newcommand{\bbN}{\mathbb{N}}
\newcommand{\bbP}{\mathbb{P}}
\newcommand{\bbR}{\mathbb{R}}
\newcommand{\bbS}{\mathbb{S}}


%\DeclareMathOperator{\cov}{Cov}
%\DeclareMathOperator{\Hess}{Hess}
%\DeclareMathOperator{\dist}{dist}
%\DeclareMathOperator{\osc}{osc}
%\DeclareMathOperator{\Id}{Id}
%\DeclareMathOperator{\diag}{diag}
%\DeclareMathOperator{\supp}{supp}
%\DeclareMathOperator{\PI}{PI}
%\DeclareMathOperator{\Ent}{Ent}




\let\a=\alpha
\let\b=\beta
\let\d=\delta
\let\e=\varepsilon
\let\g=\gamma
\let\l=\lambda
\let\o=\omega
\let\r=\rho
\let\t=\tau
\let\z=\zeta
\let\G=\Gamma
\let\O=\Omega





%%%%%%%%%%%%%%%%%%%%%%%%%%%%%%%%%%%%%%%%%%%%%%%%%%%%%%%%%%%%%%%%%%%%%%%%%%%%%%%%%%%%%%%%%%%%%%%%%%%
\graphicspath{{./figures/}}

\newcommand*{\mk}{\mkern -1mu}
\newcommand*{\Mk}{\mkern -2mu}
\newcommand*{\mK}{\mkern 1mu}
\newcommand*{\MK}{\mkern 2mu}

\hypersetup{urlcolor=purple, linkcolor=blue, citecolor=red}

\newcommand*{\romanenumi}{\renewcommand*{\theenumi}{\roman{enumi}}}
\newcommand*{\Romanenumi}{\renewcommand*{\theenumi}{\Roman{enumi}}}
\newcommand*{\alphenumi}{\renewcommand*{\theenumi}{\alph{enumi}}}
\newcommand*{\Alphenumi}{\renewcommand*{\theenumi}{\Alph{enumi}}}
\let\oldtilde\tilde
\renewcommand*{\tilde}[1]{\mathchoice{\widetilde{#1}}{\widetilde{#1}}{\oldtilde{#1}}{\oldtilde{#1}}}
\let\oldhat\hat
\renewcommand*{\hat}[1]{\mathchoice{\widehat{#1}}{\widehat{#1}}{\oldhat{#1}}{\oldhat{#1}}}
\let\oldforall\forall
\renewcommand*{\forall}{\mathrel{\oldforall}}

\usepackage{enumerate}
\newtheorem{theo}{Theorem}[section]
\newtheorem{lemm}[theo]{Lemma}
\newtheorem{prop}[theo]{Proposition}
\newtheorem{coro}[theo]{Corollary}
\theoremstyle{definition}
\newtheorem{defi}[theo]{Definition}
\theoremstyle{plain}
\newtheorem{rema}[theo]{Remark}
\numberwithin{equation}{section}
\title[Nonlinear recombinations and generalized random transpositions]{Nonlinear recombinations and generalized random transpositions: original Theorem 1.8}
\author{Pietro Caputo and Daniel Parisi}
\date{Source: Annales Henri Lebesgue 7 (2024), 1245--1299; DOI 10.5802/ahl.219}
\begin{document}
\maketitle
\noindent\textit{Editorial scope.} The counted unit is original Theorem 1.8 (particle-number-uniform entropy production for the generalized random transpositions), source0.tex lines 375--389, together with its complete author proof, source0.tex lines 2078--2136. The uncounted core retained verbatim is the nonlinear recombination model with the two crossover measures, the particle system with its generator and admissible density space, the generalized random transposition setting, Definition 1.7 of strict separation, the Dirichlet form and entropy production constants (1.12)--(1.14), and source section 4 with its notation, the weighted Shearer inequality, the base-case transposition lemma, the permutation entropies and Lemmas 4.3--4.4 with their complete proofs. Excluded and not used to inflate the difficulty: the local central limit theorem and chaos development, uniform in time propagation of chaos, the proof of Proposition 1.11, Theorems 1.12--1.13 and their proofs, and the appendix. All mathematics is native editable text with no image dependency; published equation, theorem, subsection and citation numbers are retained, and cross-references to material outside this excerpt are given as links to the published article. Printed slips are kept literal without mathematical repair: the lowercase opening ``for all $N\geq 2$'' and the unused ``for any $i=1,\,\dots,\,n$'' in the base-case Lemma 4.2, and the bracket-free conditional entropy shorthand $\mu[\ent_Vf]$ of Lemma 4.4 reused as $\mu[\ent_{A^c}f]$ in the uniform crossover computation.
\section{Introduction}\label{sec:1}

\subsection{The nonlinear equation}\label{setup}
Let $\Omega = \prod_{i=1}^nX_i$ be the set of $n$-vectors $\sigma = (\sigma_1,\,\dots,\,\sigma_n)$ where $\si_i\in X_i$, and the $X_i$ are given finite sets. We interpret $\si$ as a particle. Thus, a particle is a string of $n$ characters each taken from a finite alphabet. A basic example is obtained by taking $\Omega=\{0,1\}^n$. Without loss of generality we will assume that each $X_i$ has the form $X_i :=\{0,1,2,\,\dots,\,q_i\}$, for some $q_i \in \bbN$. Given a subset $A \subset [n],$ $[n]=\{1,\,\dots,\,n\}$, and $\sigma \in \Omega,$ $\sigma_A$ denotes the $A-$component of $\sigma,$ that is the string $(\sigma_i, i \in A).$ If $(\sigma, \eta) \in \Omega \times \Omega$ is a pair of particles, the recombination at $A$ consists in exchanging the $A$-component of $\sigma$ with the $A$-component of $\eta.$ This defines the map
\begin{align*}
(\sigma,\eta) \mapsto (\eta_A\sigma_{A^c},\sigma_A\eta_{A^c}),
\end{align*}
where $\eta_A\sigma_{A^c}$ denotes the element of $\Omega$ with entries $\eta_i$ for $i \in A$ and $\sigma_i$ for $i \in A^c=[n]\setminus A.$ Let $\cP(\O)$ denote the set of probability measures on $\O$. If the original pair $(\si,\eta)$ is obtained by sampling independently from $p\in\cP(\O)$, then the new particle $\eta_A\si_{A^c}$ is distributed according to $p_A\otimes p_{A^c}$, the product of the two marginals of $p$ on the $A$ and $A^c$ components respectively. By choosing the set $A\subset[n]$ according to some distribution $\nu$, one obtains the quadratic collision kernel
\begin{align*}
p\mapsto Q(p) = \sum_{A \subset [n]}\nu(A)\,p_{A}\otimes p_{A^c}.
\end{align*}
The nonlinear evolution is defined by the dynamical system $\dot{p}_t = Q(p_t) - p_t$, that is
\begin{align}\label{boltzeq}
\frac{d}{dt}\,p_t = \sum_{A \subset [n]}\nu(A)\left(p_{t,A}\otimes p_{t,A^c} - p_t\right)
\end{align}
with the initial condition $p_0=p\in\cP(\O)$. Here $p_t\in\cP(\O)$ is the distribution of the particle at time $t$, $p_{t,A}$ denotes its marginal on $A$ and $\nu$ is a given probability measure over the subsets of $[n]$. The study of this model starts with the pioneering work of Geiringer~\cite{Gei}; see also~\cite{Martinez,Sinetal,Sinetal2,Baakeetal} for more recent accounts. It is well known that the Cauchy problem associated to~\eqref{boltzeq} has a unique solution for every initial distribution $p\in\cP(\O)$; see e.g.~\cite{Baakeetal}. Moreover, it is not difficult to see that the evolution preserves the single site marginals, that is $p_{t,i}=p_{i}$ for all $t\geq 0$ and for all $i\in[n]$. We say that the recombination measure $\nu$ is \emph{separating} if for any $i, j \in [n]$ there is a positive probability that the random set $A$ with distribution $\nu$ separates $i$ and $j$. Equivalently, if $r(\nu)<1$ where
\begin{equation}\label{def:rnu}
r(\nu):= \max_{i\,<\,j\,\in\,[n]}\,\nu\left(\{i,j\}\subset A \text{ or }\{i,j\}\subset A^c\right)
\end{equation}
denotes the maximum over $i<j$ of the probability that $A$ does not separate $i,j$. It is a classical fact that, under the assumption that $\nu$ is separating, the system converges to the stationary state given by the product of the marginals of the initial state $p$; namely, if $\pi_i=p_{i}$ denotes the marginal of $p$ at site $i$, then
\begin{equation}\label{station}
\pi=\otimes_{i=1}^n\,\pi_i
\end{equation}
is the equilibrium distribution and one has the convergence $p_t\to \pi$, $t\to\infty$, which can be interpreted as the effect of repeated fragmentations of the initial state; see e.g.~\cite{Sinetal2}. Some of our results will hold for arbitrary separating $\nu$. In some other cases we consider a slightly stronger assumption on $\nu$. Two examples to which all our results apply are the following distributions $\nu$, which are commonly considered in the genetic recombination literature. We refer to~\cite{entprod,Sinetal2} for more examples.
\begin{enumerate}
\item \emph{Uniform crossover}: $\nu(A)=\frac1{2^n}$, for all $A\subset[n]$;
\item \emph{One-point crossover}: $\nu(A)=\tfrac1{n+1}\sum_{i=0}^n\ind_{A=J_i} $, where $J_0=\emptyset$, $J_i=\{1,\,\dots,\,i\}$, $i\geq 1$.
\end{enumerate}

\setcounter{subsection}{1}
\setcounter{equation}{3}
\subsection{The particle system}

Suppose there are $N$ ``particles'', described by variables $\eta(j)\in\O$, $j=1,\,\dots,\,N$. That is, each particle is a single string from $\O$ and $\eta_i(j)$ denotes the content of the $j^{\rm th}$ particle at site $i\in[n]$. We may picture $\eta\in\O^N$ as a $N\times n$ matrix such that each row is a particle with $n$ entries, and for each $i\in[n]$, the $i^{\rm th}$ column $\eta_i$ represents the content of site $i$ for different particles.

Notice that $N$ and $n$ play two very different roles here. The number $N$ of particles will eventually be taken to $+\infty$ to recover the non linear mean field limit, in accordance with the general Kac program. The number $n$ should be thought as a fixed, possibly large quantity describing the size of a single particle space.

The Markov process is given by the following random pair-exchange process. Pairs of particles $\{j,l\}$, $1\leq j<l\leq N$ are chosen independently according to a Poisson clock process with rate $1/N$. When the pair $\{j,l\}$ ``rings'', then a set $A\subset[n]$ is chosen with probability $\nu(A)$ and the recombination
\begin{align}\label{eqtran}
(\eta(l),\eta(j)) \mapsto \big(\eta_A(j)\eta_{A^c}(l),\eta_A(l)\eta_{A^c}(j)\big)
\end{align}
is performed, that is the $A$-content is exchanged between particle $j$ and particle $l$. For all $j,l \in [N], A \subset [n],$ for all $\eta \in \Omega^N$, we write $\eta^{j,l,A}$ for the new configuration $\eta' \in \Omega$ defined by
\begin{align}\label{def:etajl}
\eta'(k)=\eta(k), \;\forall k \neq j,l;\qquad \eta'(l)=\eta_{A}(j)\eta_{A^c}(l),\qquad \eta'(j)=\eta_A(l)\eta_{A^c}(j).
\end{align}
With this notation the pair configuration in the right hand side of~\eqref{eqtran} is $(\eta^{j,l,A}(l),\linebreak\eta^{j,l,A}(j))$. Define also $f^{j,l,A}(\eta)=f(\eta^{j,l,A})$ for all $f:\Omega^N \rightarrow \bbR$. Then, the infinitesimal generator of the continuous time Markov process is given by
\begin{align}\label{operator}
\cL_N f = \frac{1}{N}\sum_{1\,\leq\,j\,<\,l\,\leq\,N}\;\sum_{A\,\subset\,[n]}\nu(A)\left(f^{j,l,A}-f\right), \qquad f:\Omega^N \rightarrow \bbR.
\end{align}
Any product measure $\mu_N$ on $\Omega^N$ of the form $\mu_N = \mu^{\otimes N},$ where $\mu$ is itself a product measure $\mu=\mu_1\otimes\cdots\otimes\mu_n$ on $\Omega=\prod_{i=1}^nX_i,$ defines a reversible measure for the generator $\cL_N$. Indeed, for such a measure one has the symmetry $\mu_N(\eta^{j,l,A})=\mu_N(\eta)$ for all $\eta\in\O^N$ and therefore $\cL_N$ is self-adjoint in $L^2(\mu_N)$. The process is not irreducible in the state space $\Omega^N$ since the content at a site for one particle is always exchanged with the content at the same site for another particle, and thus the number of particles with a given element $x \in X_i$ at a given site $i \in [n]$ is constant in time. To obtain an irreducible process one must fix the densities $\rho_{i,x},$ $i \in [n],x\in X_i$ defined by
\begin{align*}
\rho_{i,x} = \frac{1}{N}\sum_{j=1}^N\ind(\eta_i(j) = x).
\end{align*}
We call $\rho = (\rho_{i,x})$ the corresponding vector. Given a density vector $\rho$, consider the space
\begin{multline}\label{omegarhosets}
\Omega_{\rho} := \\
\Big\{(\eta(1),\,\dots,\,\eta(N)) \in \Om^N \!\!:\, \;
\textstyle{\rho_{i,x} = \frac{1}{N}\sum_{j=1}^N \ind(\eta_i(j) = x)},\;\, \forall i \in [n], \ x \in X_i\Big\}.
\end{multline}
The set $\Om_{\rho}$ is well defined and non-empty for every vector $\r=\r_N$ such that $\r_{i,x}\in[0,1]$, $\sum_{x\,\in\,X_i}\r_{i,x}=1$ for all $i\in[n]$, and such that $N\r_{i,x}$ is an integer for all $i,x$. When this holds we say that $\r_N$ is an \emph{admissible sequence}. Under suitable assumptions on the recombination measure $\nu$, see Definition~\ref{def:strict}, for any given admissible $\r_N$, the Markov process with state space $\O_{\r_N}$ and generator $\cL_N$ is irreducible and converges to the uniform distribution on $\O_{\r_N}$. We will be interested in quantitative statements about this convergence.

\setcounter{subsection}{4}
\subsection{Entropy production for generalized random transpositions}

Our second main result is about quantitative estimates on the decay to equilibrium for the particle system introduced above. The key feature is that these estimates hold \emph{uniformly in the number of particles} $N$. We shall actually derive such estimates in the context of the \emph{generalized random transposition process} defined as follows.

Let $\cS_{N,n}=S_N^n$ denote the $n$-fold product of the symmetric group $S_N$ of the permutations of $[N]=\{1,\,\dots,\,N\}$. Then $\eta\in \cS_{N,n}$ is a matrix $\eta_i(j)$, $i\in[n]$, $j\in[N]$, where each $\eta(j)=(\eta_1(j),\,\dots,\,\eta_n(j))\in[N]^n$ is seen as a particle, and each $\eta_i=(\eta_i(1),\,\dots,\,\eta_i(N))\in S_N$ is a permutation of $[N]$. Note that $\cS_{N,n} = \O_{\r}$ where $\O_{\r}$ is defined in~\eqref{omegarhosets} when we take the extreme case $q_i\equiv N-1$, $\r_{i,x}\equiv 1/N$ for all $i=1,\,\dots,\,n$.

The generalized random transposition (GRT) process is defined as the process generated by the operator $\cL_N$ in this setup, namely the GRT process is the continuous time Markov process with state space $\cS_{N,n}$ described as follows: every pair of particles $\{\eta(l),\eta(j)\}$ collides with rate $1/N$ independently, and when a collision occurs, a new set $A$ is sampled according to $\nu$ and the $A$-content of $\eta(l),\eta(j)$ is exchanged.

This setting is convenient for proving functional inequalities since by restricting to classes of functions with suitable symmetries we then recover all possible cases of processes on $\O_{\r_N}$ with generator $\cL_N$, for all admissible $\r_N$, see Remark~\href{https://ahl.centre-mersenne.org/item/AHL_2024__7__1245_0/}{1.9} below for more details. As an example, consider the case $q_i\equiv 1$ and suppose that $N(\r_N)_{i,1} = N_i$ for some positive integers $N_i$, $i=1,\,\dots,\, n$. Here the process generated by $\cL_N$ can be seen as the GRT process restricted to functions $f:\cS_{N,n}\mapsto\bbR$ such that, for each $i$, $f$ only depends on $(\eta_i(1),\,\dots,\,\eta_i(N))$ through the unordered set $\{\eta_i(1),\,\dots,\,\eta_i(N_i)\}$. This can be seen as a generalized Bernoulli--Laplace process~\cite{DS}. If no confusion arises we continue to write $\cL_N$ for the generator of the GRT.


We remark that when $n=1$, GRT is just the usual random transposition process~\cite{DiacShash}, and that when $\nu$ gives positive weight only to $A\subset [n]$ such that $|A|=1$, it describes $n$ independent random transposition processes. However, in the general case, the recombination measure $\nu$ dynamically couples the permutations and the GRT becomes a nontrivial generalization of the standard random transpositions.


In order to guarantee the irreducibility of the GRT process, we make the following assumption on the recombination measure $\nu$, which is easily seen to be stronger than the separation assumption $r(\nu)<1$; see also Remark~\href{https://ahl.centre-mersenne.org/item/AHL_2024__7__1245_0/}{1.10}.

\setcounter{theo}{6}
\begin{defi}\label{def:strict}
We say that $\nu$ is \emph{strictly separating} if for all $i\in[n]$ there exists $A\subset [n]$ such that $i\in A$ and such that both $A$ and $A\setminus\{i\}$ have positive $\nu$-probability.
\end{defi}

\setcounter{equation}{11}
Note that the uniform crossover and the one-point crossover are both strictly separating. Let $\pi_N$ denote the uniform distribution on $\cS_{N,n}$. The GRT process is reversible with respect to $\pi_N$ and if the measure $\nu$ is strictly separating, then it is also irreducible, and any initial distribution converges to $\pi_N$ as $t\to\infty$. To quantify this statement we consider the Dirichlet form of the GRT, defined by

\begin{multline}\label{operatorG}
\cE_{N,n}(f,g)\\
=\frac{1}{2N}\sum_{1\,\leq\,j\,<\,l\,\leq\,N}\sum_{A\,\subset\,[n]}\nu(A)
\sum_{\eta\,\in\,\cS_{N,n}}\pi_N(\eta)\left(f\left(\eta^{j,l,A}\right)-f(\eta)\right)\left(g\left(\eta^{j,l,A}\right)-g(\eta)\right),
\end{multline}
where $f,g:\cS_{N,n} \mapsto\bbR$. The \emph{entropy production rate} is measured by the constant
\begin{align}\label{deltaoperatorG}
\a(N,n) = \inf_{f\,>\,0} \frac{\cE_{N,n}(f,\log f)}{\ent(f)},
\end{align}
where the infimum is over $f:\cS_{N,n} \mapsto\bbR_+$ such that $\ent(f)\neq 0$ and
\[
\ent(f) = \pi_N(f\log f)-\pi_N(f)\log\pi_N(f)
\]
is the entropy of $f$ w.r.t.\ $\pi_N$.
Equivalently, $\a(N,n)$ is the best constant $\a$ such that the inequality
\begin{align}\label{gapoperatorG}
\ent\left(e^{t\cL_N}f\right)\leq e^{-\a t}\ent(f)
\end{align}
holds for all functions $f>0$; see e.g.~\cite{BT,DS}. Note that when $\pi_N(f)=1$ then, for all $t\geq 0$, $\ent(e^{t\cL_N}f)$ coincides with the relative entropy $H_N(\mu_{N,t}\tc\pi_N)$ where $\mu_{N,t} =(e^{t\cL_N}f)\pi_N$.

We also consider the entropy production rate restricted to the set of symmetric functions defined as follows. Let $\bbS$ denote the set of $f:\O^N\mapsto \bbR$ such that
\[
f(\eta)=\frac1{N!}\sum_{\t\,\in\,S_N} f(\t\circ\eta),
\]
where the sum runs over all permutations $\t\in S_N$ and $\t\circ\eta$ denotes the configuration with particles exchanged according to $\t$, that is $\t\circ\eta(j) = \eta(\t(j))$. From the point of view of Kac's program~\cite{KAC}, $\bbS$ is the relevant space of observables in the particle system. We call $\a_{\bbS}(N,n)$ the constant defined as in~\eqref{deltaoperatorG}, with the infimum restricted to positive functions $f\in\bbS$.

\setcounter{theo}{7}
\setcounter{equation}{14}
\begin{theo}
\label{th:ent}
Fix $n\in\bbN$ and assume that $\nu$ is strictly separating. Then there exists $\a(\nu)>0$ such that for any $N\in\bbN$, $N\geq 2$,
\begin{equation}\label{entao}
\a(N,n)\geq \a(\nu).
\end{equation}
Moreover, if $\nu$ is the one-point crossover, then
\begin{equation}\label{enta}
\a(N,n)\geq \frac{1}{4(n+1)}\,,
\end{equation}
and if $\nu$ is the uniform crossover, then
\begin{equation}\label{entas}
\a(N,n)\geq \frac{1}{4n}\,,\qquad \a_\bbS(N,n)\geq \frac1{2(n+2)}\,.
\end{equation}
\end{theo}

\setcounter{section}{3}
\section{Relative entropy decay}\label{sec:entdec}

The goal of this section is to prove Theorem~\ref{th:ent}, Proposition~\href{https://ahl.centre-mersenne.org/item/AHL_2024__7__1245_0/}{1.11}, and Theorem~\href{https://ahl.centre-mersenne.org/item/AHL_2024__7__1245_0/}{1.12}. We start by setting up a convenient notation and by gathering some preliminary ingredients of the proof. To simplify our notation we write $\mu=\pi_N$ for the uniform distribution over $\cS_{N,n}$. For any $A\subset [n]$, $f:\cS_{N,n}\mapsto\bbR$, we are going to use the notation

\[
\mu_A f =\mu\left(f\,\middle|\,  \eta_{A^c}\right),
\]
where $\mu(\cdot\tc \eta_{A^c})$ denotes the conditional expectation given the variables
\[
\eta_{A^c}:=\{\eta_i(j), \;i\in A^c, \;j=1,\,\dots,\,N\}.
\]
The function $\mu_Af$ thus depends on $\eta$ only through the variables $\eta_{A^c}$. When $A=[n]$ we have the global expectation $\mu_{[n]} f = \mu(f)$. With this notation, $\mu_A$ is the orthogonal projection, in $L^2(\cS_{N,n},\mu)$, onto the space of functions that do not depend on the $A$-component of $\eta\in\cS_{N,n}$.
Notice the relations $\mu(f)=\mu(\mu_A f)$, and $\mu_B(f)=\mu_B(\mu_A f)$, for all $A\subset B\subset [n]$. We also use the notation $\ent_A(f)$ for the entropy of $f:\cS_{N,n}\mapsto\bbR_+$ with respect to $\mu_A$, that is
\begin{equation}\label{entA}
\ent_A(f) = \mu_A\left[f\log(f/\mu_A f)\right]\,.
\end{equation}
Thus $\ent_A(f)$ is a function that depends on $\eta$ through the variables $\eta_{A^c}$ only. Its expectation satisfies
\begin{equation}\label{entA2}
\mu\left[\ent_A(f)\right] = \ent(f) - \ent[\mu_Af]\,,
\end{equation}
where $\ent(f) = \ent_{[n]}(f) = \mu\left[f\log(f/\mu (f))\right]$. In our setting the well known Shearer inequality can be formulated as follows.

\begin{lemm}\label{lem:shearer}
For any distribution $\nu$ on subsets of $[n]$, for any $f:\cS_{N,n}\mapsto\bbR_+$,
\begin{equation}\label{shearer}
\sum_{A\subset [n]}\nu(A) \,\mu\left[\ent_A(f)\right] \geq \g(\nu)\, \ent(f)\,,
\end{equation}
where $\g(\nu) = \min_{i\,\in\,[n]}\;\sum_{A:\,A\,\ni\,i}\nu(A)$.
\end{lemm}

\begin{proof}
Reasoning as in~\cite[Proposition~4.3]{entprod}, the estimate is a consequence of the product structure of $\mu$ and the weighted version of the classical Shearer inequality for Shannon entropy. We refer e.g.\ to~\cite{Virag} for a proof of the latter.
\end{proof}

When $A=\{i\}$, a single site, we write $\mu_i$ for $\mu_{\{i\}}$ and $\ent_i$ for $\ent_{\{i\}}$. An immediate consequence of Lemma~\ref{lem:shearer} is the following well known tensorization property:
\begin{equation}\label{tensor}
\ent(f) \leq\sum_{i=1}^n \mu\left[\ent_{i}(f) \right].
\end{equation}



\subsection{The case \texorpdfstring{$n=1$}{n=1}}

A key ingredient of our proof is the control of the base case $n=1$. Here the problem reduces to standard random transpositions and one can use a well known bound that was first derived in~\cite{Goel,GQ}. In our setting it can be summarized as follows, see~\cite[Theorem~1]{GQ} or~\cite[Corollary~3.1]{Goel} for a proof. Note that for $n=1$ we have $\cS_{N,n} =S_N$.

\begin{lemm}
\label{GQlemma} for all $N\geq 2$, for all $g:S_N\mapsto\bbR_+$, for any $i=1,\,\dots,\,n$,
\begin{align*}
\sum_{\t\,\in\,S_N} &g(\t)\log(g(\t)/\bar g)
\leq \frac1N\sum_{j<\ell}
\sum_{\t\,\in\,S_N}
\left(g(\t^{j,\ell}) - g(\t)\right)\log \frac{g(\t^{j,\ell})}{g(\t)},
\end{align*}
where $\bar g = \frac1{N!}\sum_{\t\,\in\,S_N}g(\t)$, and $\t^{j,\ell}$ denotes $\t$ composed with the transposition at $\{j,\ell\}$.
\end{lemm}

The proof of Theorem~\ref{th:ent} is based on Lemma~\ref{lem:shearer}, Lemma~\ref{GQlemma} and the following analysis of the entropy associated to partial random permutations of the particle configuration, which is the main novelty in this section.

\subsection{Permutation entropies}

For any $A\subset [n]$, $g:\cS_{N,n}\mapsto\bbR$, define $P_Ag:\cS_{N,n}\mapsto \bbR$, as
\begin{align}\label{defpA}
P_Ag(\eta) = \frac1{N!}\sum_{\t\,\in\,S_N}g((\t\eta)_A\,\eta_{A^c}),
\end{align}
where $\t\eta:=\t\circ\eta$ denotes the element of $\cS_{N,n}$ obtained from $\eta$ by permuting the particle labels according to $\t$:
\[
(\t\eta)(j) = \eta(\t(j)).
\]
The linear operator $P_A$ is the orthogonal projection, in $L^2(\cS_{N,n},\mu)$, onto the space of functions that are symmetric w.r.t.\ permutations restricted to the subset $A$. When $A=\{i\}$ we write $P_{\{i\}}=P_i$, and note that $P_{i}g = \mu(g\tc \eta_{\{i\}^c})=\mu_i g$ for all $i$. Note also that $P_A,P_B$ do not commute for general $A,B\subset[n]$, but if $A\cap B=\emptyset$ then $P_AP_B=P_BP_A$. Moreover, one easily checks that if $A\subset B\subset [n]$, then
\begin{align}\label{commPA}
P_AP_B = P_{B\setminus A}P_A = P_AP_{B\setminus A}=P_BP_A,\qquad \;A\subset B.
\end{align}
Also, observe that the orthogonal projection $\mu_A$ defined above satisfies
\begin{align}\label{muAPA}
\mu_A = \prod_{i\in A}P_i.
\end{align}
Notice that $P_{[n]}g = g$ iff $g\in\bbS$, and that in this case
\begin{align}\label{enta+7}
P_Ag = P_{A^c}g = P_AP_{A^c}g=P_{A^c}P_{A}g\,,\qquad g\in\bbS.
\end{align}
For any fixed $f\geq 0$, and $A\subset[n]$, define
\begin{align}\label{phia}
\phi(A;f) = \mu\left[f\log(f/P_Af)
\right].
\end{align}
The function $\phi$ is a suitable conditional entropy of $f$ on $A$. Namely, $\phi(A;f)$ is the expected value with respect to $\mu$ of the entropy
\[
P_A\left[f\log(f/P_Af)\right].
\]
In particular, $\phi(A;f) \geq 0$. We call $\phi(A;f) $ the \emph{permutation entropy} of $f$ on $A$. When there is no risk of confusion we write simply $\phi(A)$ for $\phi(A;f)$. In general, one has

\begin{lemm}
\label{propphiA}
Fix a function $f\geq 0$. For any $i\in[n]$, $\phi(\{i\}) = \mu[\ent_i(f)]$, and for all $A\subset[n]$:
\begin{equation}\label{phia2}
\phi(A) = \mu\left[\ent_A(f) \right] - \mu\left[\ent_A(P_Af) \right] \,.
\end{equation}
Moreover, for all $A\subset[n]$
\begin{equation}\label{phiaGQ}
\phi(A) \leq 2\cE_A(f,\log f),
\end{equation}
where
\begin{align}\label{enta+5}
\cE_A(f,\log f)=\frac12\frac1N\sum_{j\,<\,\ell}\mu\left[\left(f^{j,\ell,A} - f\right)
\log \frac{f^{j,\ell,A}}{f}
\right].
\end{align}
\end{lemm}


\begin{proof}
If $A=\{i\}$, then $P_A f = \mu_i f$ and therefore $\phi(\{i\}) = \mu\left[\ent_i(f)\right]$. In general, for any $A$,
\begin{align*}
\mu\left[\ent_A(f) \right]&= \mu\left[f\log(f/\mu_Af)\right] \\
& =\mu\left[f\log(f/P_Af)\right] + \mu\left[f\log(P_Af/\mu_Af)\right].
\end{align*}
The second term above satisfies
\[
\mu\left[f\log(P_Af/\mu_Af)\right]= \mu\left[P_A f\log(P_Af/\mu_A P_Af)\right] = \mu\left[\ent_A(P_Af) \right] \,,
\]
where we have used that $\mu_BP_A=\mu_B$ for any $A\subset B\subset[n]$ since $P_A$ is an orthogonal projection. This proves~\eqref{phia2}.

To prove~\eqref{phiaGQ}, observe that for any $A\subset [n]$, $f\geq 0$, any fixed $\eta\in\cS_{N,n}$, taking $g(\t):=f((\t\circ\eta)_A\eta_{A^c})$, one has $\bar g = P_Af$, and therefore by Lemma~\ref{GQlemma},
\begin{multline*}
P_A\left[f\log(f/P_Af)\right] (\eta)=\frac1{N!}\sum_{\t\,\in\,S_N} g(\t)\log(g(\t)/\bar g)\\
\begin{aligned}
&\leq \frac1N\;\sum_{j\,<\,\ell}\frac1{N!}\sum_{\t\,\in\,S_N}\left(f\left((\t^{j,\ell}\eta)_A\eta_{A^c}\right) - f\left((\t\eta)_A\eta_{A^c}\right)\right)\log \frac{f\left((\t^{j,\ell}\eta)_A\eta_{A^c}\right)}{f\left((\t\eta)_A\eta_{A^c}\right)}
\\&=
\frac1N\sum_{j\,<\,\ell}P_A\left[\left(f^{j,\ell,A} - f\right)\log \frac{f^{j,\ell,A}}f
\right](\eta).
\end{aligned}
\end{multline*}
Taking the expectation and using $\mu P_A=\mu$ one obtains~\eqref{phiaGQ}.
\end{proof}

Next, we compare the permutation entropy $\phi(A)$ with the usual entropy $\mu[\ent_A(f) ]$. The previous lemma in particular shows that $\phi(A)\leq \mu[\ent_A(f) ]$. The next lemma allows us to give a bound in the opposite direction. Notice that such a bound needs some care since if e.g.\ $f\in\bbS$ is not a constant then $\phi([n];f)=0$ while $\mu[\ent_{[n]}(f) ]=\ent(f)\neq 0$.

\begin{lemm}\label{keyphi}
Fix $A\subset [n]$ such that $0\leq |A|\leq n-1$, and $V\subset A^c$. Then,
\begin{equation}\label{phia5}
\phi(A) + \sum_{i\,\in\,V}\phi(A\cup\{i\})\geq \mu\left[\ent_{V}f \right].
\end{equation}
\end{lemm}

\begin{proof}
Write $V=\{x_1,\,\dots,\,x_k\}\subset [n]$, $k=|V|$, and define $A_0=A$, and $A_i=A\cup\{x_i\}$, $i=1,\,\dots,\,k$. In general $P_{A_i}$ and $P_{A_j}$ do not commute for $i,j=1,\,\dots,\,k$, but using~\eqref{commPA} one has $P_{A_0}P_{A_i} = P_{A_0}P_{x_i}= P_A\mu_{x_i}$ and for any $\ell=1,\,\dots,\,k$,
\[
P^\ell:=P_{A_0}P_{A_1}\cdots P_{A_\ell} = P_A\prod_{i=1}^\ell \mu_{x_i}= \left(\prod_{i=1}^\ell \mu_{x_i} \right)P_A =\prod_{i=0}^\ell P_{A_i}\,,
\]
where the last identity holds regardless of the order of the multiplications. In other words, the operators $P_{A_i}$ commute thanks to the presence of $P_{A_0}$. In particular, $P^k=P_A\mu_V=\mu_VP_A$.

Consider the entropy
\begin{align}\label{eq:summa}
\mu\left[f\log\frac{f}{P^kf}\right] = \mu\left[f\log\frac{f}{\mu_Vf}\right] + \mu\left[f\log\frac{\mu_Vf}{P_A\mu_Vf}\right].
\end{align}
The first term is $\mu\left[\ent_Vf\right]$. The second term satisfies
\begin{align}\label{eq:summaa}
\mu\left[f\log\frac{\mu_Vf}{P_A\mu_Vf}\right] = \mu\left[\mu_Vf\log\frac{\mu_Vf}{\mu_VP_Af}\right] = \phi(A;\mu_Vf)\geq 0.
\end{align}
On the other hand,
\[
\log\frac{f}{P^kf}=\log\frac{f}{P_{A_0} f} + \log\frac{P_{A_0}f}{P_{A_0}P_{A_1}f}+\cdots+\log\frac{P_{A_0}P_{A_1}\cdots P_{A_{k-1}} f}{P_{A_0}P_{A_1}\cdots P_{A_k} f},
\]
and therefore
\begin{align}\label{eq:summ}
\mu\left[f\log\frac{f}{P^kf}\right] = \phi(A;f)+\sum_{\ell=0}^{k-1}\mu\left[f\log\frac{P^\ell f}{P^{\ell+1}f}\right].
\end{align}
Next, we show that
\begin{equation}\label{phia05}
\mu\left[f\log\frac{P^\ell f}{P^{\ell+1}f}\right]
\leq \phi\left(A_{\ell+1};f\right),
\end{equation}
for all $\ell=0,\,\dots,\,k-1$. Combined with~\eqref{eq:summa}-\eqref{eq:summ}, this implies the desired conclusion:
\begin{equation}\label{phia5x}
\phi(A) + \sum_{i\,\in\,V}\phi(A\cup\{i\})\geq \mu\left[\ent_{V}f \right] + \phi(A;\mu_V f)\geq \mu\left[\ent_{V}f \right].
\end{equation}
It remains to prove~\eqref{phia05}. We first observe that
\begin{equation}\label{phia5xx}
\mu\left[f\log\frac{P^\ell f}{P^{\ell+1}f}\right]=\mu\left[P^\ell f\log\frac{P^\ell f}{P^{\ell+1}f}\right]=\phi\left(A_{\ell+1};P^\ell f\right).
\end{equation}
The well known variational principle for the relative entropy implies that for any $B\subset [n]$
\[
P_B\left(f\log\frac{f}{P_Bf}\right)\geq P_B(fg)\,,
\]
for any function $g$ such that $P_B(e^g)=1$. Choosing $g= \log\frac{P^\ell f}{P_BP^\ell f}$ shows that
\[
P_B\left(f\log\frac{f}{P_Bf}\right)\geq P_B\left(f\log\frac{P^\ell f}{P_BP^\ell f}\right)\,.
\]
Taking the expectation one finds
\[
\phi(B;f)\geq\mu\left[f\log\frac{P^\ell f}{P_BP^\ell f}
\right]\,.
\]
If $P_B,P^\ell $ commute, then
\[
\mu\left[f\log\frac{P^\ell f}{P_BP^\ell f}
\right]= \mu\left[f\log\frac{P^\ell f}{P^\ell P_Bf}
\right]=\mu\left[P^\ell f\log\frac{P^\ell f}{P^\ell P_Bf}
\right] = \phi\left(B;P^\ell f\right).
\]
Therefore,
\[
\phi(B;f)\geq\phi\left(B;P^\ell f\right),
\]
whenever $P_B,P^\ell $ commute. Taking $B=A_{\ell+1}$, and using the fact that $P_{A_{\ell+1}}$ and $P^{\ell}$ commute, one obtains $\phi\left(A_{\ell+1};P^\ell f\right)\leq \phi\left(A_{\ell+1}; f\right)$. Together with~\eqref{phia5xx}, this implies~\eqref{phia05}.
\end{proof}

\subsection{Proof of Theorem~\ref{th:ent}}

From the strict separation assumption, it follows that for some $\d(\nu)>0$,
\[
\cE_{N,n}(f,\log f)
\geq \frac{\d(\nu)}n\sum_{i=1}^n\left(\cE_{A_i}(f,\log f)+\cE_{A_i\setminus\{i\}}(f,\log f)\right),
\]
where, for every $i$, $A_i\subset[n]$ is such that $A_i\ni i$ and $\min\{\nu(A_i),\nu(A_i\setminus\{i\})\}\geq \d(\nu)$. Therefore, from Lemma~\ref{propphiA},
\begin{equation}\label{phia5a1}
\cE_{N,n}(f,\log f)
\geq \frac{\d(\nu)}{2n}\sum_{i=1}^n\left(\phi(A_i) +\phi(A_i\setminus\{i\})\right).
\end{equation}
Lemma~\ref{keyphi} then implies
\begin{equation}\label{phia5a2}
\cE_{N,n}(f,\log f)
\geq \frac{\d(\nu)}{2n}\sum_{i=1}^n\mu\left[\ent_i f\right].
\end{equation}
Using~\eqref{tensor} it follows that $\a(N,n)\geq \d(\nu)/{2n}$. This proves the bound~\eqref{entao} with $\a(\nu)=\d(\nu)/{2n}$.

To prove the lower bound for one-point crossover, notice that the above argument can be repeated but this time we can directly estimate
\[
\cE_{N,n}(f,\log f)\geq \frac14\frac1{n+1}\sum_{i=1}^n (\phi(J_i)+\phi(J_{i-1}))\geq
\frac1{4(n+1)}\sum_{i=1}^n\mu\left[\ent_{i}f \right],
\]
where $J_0=\emptyset$, and $J_i=\{1,\,\dots,\,i\}$, $i\geq 1$. The lower bound $\a(N,n)\geq 1/4(n+1)$ thus follows again by the tensorization~\eqref{tensor}.

Next, we prove the lower bound for the case of uniform crossover $\nu(A)=2^{-n}$ for all $A\subset[n]$. From Lemma~\ref{propphiA}
\[
\cE_{N,n}(f,\log f)=2^{-n}\sum_A \cE_A(f,\log f)\geq 2^{-n-1}\sum_A \phi(A).
\]
By Lemma~\ref{keyphi},
\begin{multline*}
\sum_{i=1}^n\sum_{A}\phi(A)\ind(i\in A)\\
= \sum_{i=1}^n\;\sum_{A: \,|A|\leq n-1}\phi(A\cup\{i\})\ind(i\notin A)
\geq \sum_{A: \,|A|\,\leq\,n-1} \left(\mu\left[\ent_{A^c}f \right] - \phi(A)\right).
\end{multline*}
Therefore,
\begin{align}\label{enta+43}
\begin{split}
\sum_A \phi(A)&= \frac1n\sum_{i=1}^n\sum_{A}\phi(A)(\ind(i\notin A) +\ind(i\in A))
\\&\geq \frac1n\sum_{A}|A^c|\phi(A) + \frac1n\sum_{A: \,|A|\,\leq\,n-1}\left(\mu\left[\ent_{A^c}f \right] - \phi(A)\right) \\
&=\frac1n\sum_{A:\, |A|\,\leq\,n-1}(|A^c|-1)\phi(A) + \frac1n\sum_{A}\mu\left[\ent_{A^c}f \right].
\end{split}
\end{align}
In particular,
\[
2^{-n-1}\sum_A \phi(A)\geq \frac{2^{-n}}{2n}\sum_A\mu\left[\ent_{A^c}f \right].
\]
From Lemma~\ref{lem:shearer} it follows that
\[
2^{-n-1}\sum_A \phi(A)\geq \frac{1}{4n}\ent f.
\]
This proves the desired bound $\a(N,n)\geq \frac{1}{4n}$.

Finally, in the case of symmetric functions one has $\phi(A)=\phi(A^c)$ for all $A\subset[n]$, see~\eqref{enta+7}. Therefore, the previous computation now shows that
\[
\sum_A \phi(A)\geq \frac2n\sum_{A}\mu\left[\ent_{A^c}f \right] - \frac2n \sum_A \phi(A).
\]
Rearranging terms yields the bound $\a_\bbS(N,n)\geq \frac{1}{2(n+2)}$. This concludes the proof of Theorem~\ref{th:ent}.

\begin{thebibliography}{99}
\bibitem[Gei44]{Gei} H. Geiringer, ``On the probability theory of linkage in Mendelian heredity'', \emph{Ann. Math. Stat.}\textbf{15}(1944), 25--57.
\bibitem[Mar17]{Martinez} S. Mart\'inez, ``A probabilistic analysis of a discrete-time evolution in recombination'', \emph{Adv. Appl. Math.}\textbf{91}(2017), 115--136.
\bibitem[RSW92]{Sinetal} Y. Rabinovich, A. Sinclair and A. Wigderson, ``Quadratic dynamical systems'', 33rd annual symposium on Foundations of computer science (FOCS), IEEE, 1992, 304--313.
\bibitem[RRS98]{Sinetal2} Y. Rabani, Y. Rabinovich and A. Sinclair, ``A computational view of population genetics'', \emph{Random Struct. Algorithms}\textbf{12}(1998), no.4, 313--334.
\bibitem[SBB16]{Baakeetal} M. Salamat, M. Baake and E. Baake, ``The general recombination equation in continuous time and its solution'', \emph{Discrete Contin. Dyn. Syst.}\textbf{36}(2016), no.1, 63--95.
\bibitem[CAS18]{entprod} P. Caputo and A. Sinclair, ``Entropy production in nonlinear recombination models'', \emph{Bernoulli}\textbf{24}(2018), no.4B, 3246--3282.
\bibitem[DSC96]{DS} P. Diaconis and L. Saloff-Coste, ``Logarithmic Sobolev inequalities for finite Markov chains'', \emph{Ann. Appl. Probab.}\textbf{6}(1996), no.3, 695--750.
\bibitem[DS81]{DiacShash} P. Diaconis and M. Shahshahani, ``Generating a random permutation with random transpositions'', \emph{Z. Wahrscheinlichkeitstheor. Verw. Geb.}\textbf{57}(1981), 159--179.
\bibitem[BT06]{BT} S. G. Bobkov and P. Tetali, ``Modified logarithmic Sobolev inequalities in discrete settings'', \emph{J. Theor. Probab.}\textbf{19}(2006), no.2, 289--336.
\bibitem[Kac56]{KAC} M. Kac, ``Foundations of kinetic theory'', Proceedings of the Third Berkeley Symposium on Mathematical Statistics and Probability, 1954--1955, vol. III, University of California Press, 1956, 171--197.
\bibitem[CHV20]{Virag} E. Cs\'oka, V. Harangi and B. Vir\'ag, ``Entropy and expansion'', \emph{Ann. Inst. Henri Poincar\'e, Probab. Stat.}\textbf{56}(2020), no.4, 2428--2444.
\bibitem[Goe04]{Goel} S. Goel, ``Modified logarithmic Sobolev inequalities for some models of random walk'', \emph{Stochastic Processes Appl.}\textbf{114}(2004), no.1, 51--79.
\bibitem[GQ03]{GQ} F. Gao and J. Quastel, ``Exponential decay of entropy in the random transposition and Bernoulli--Laplace models'', \emph{Ann. Appl. Probab.}\textbf{13}(2003), no.4, 1591--1600.
\end{thebibliography}
\end{document}
